\documentclass[11pt,letterpaper]{article}
\usepackage{hanzo-defense}

\doctitle{Sovereign vs.\ Closed: Hanzo AI versus\\Centralized Cloud AI for Defense and\\Regulated Workloads}
\docid{HAI-DEF-2026-011}
\docversion{Version 1.0}
\docdate{2026}
\docauthors{Zach Kelling --- Hanzo Team}
\docorgs{Hanzo Industries\\research@hanzo.ai}

\begin{document}
\makecover

\begin{abstract}
We compare two structurally different ways to obtain frontier AI: a closed,
centrally hosted model accessed as a metered cloud service (typified by
Anthropic's Claude), and a sovereign, self-hostable stack with owned model
weights (Hanzo). The distinction is not primarily one of model quality---closed
frontier models are excellent---but of \emph{deployment model and control}.
For defense and regulated workloads, the binding constraints are data residency,
operation in disconnected and contested environments, supply-chain provenance,
post-quantum security, auditability, and cost at scale. A closed cloud service
cannot satisfy several of these by construction: it cannot run air-gapped, its
weights cannot be inspected or owned, and sensitive data must leave the
enclave to be processed. We argue that for these workloads the relevant choice
is between foreign open-weight models and an American open-source alternative,
and we position Hanzo as the latter---owning both the model (the Zen family)
and the cryptography (a production post-quantum stack). We give a dimension-by-dimension
comparison and a decision framework for when each model of consumption fits.
\end{abstract}

\paragraph{Executive Summary.}
This is not a claim that our model is smarter than Claude. It is a claim about
\emph{where} and \emph{how} you can run it. Claude is an excellent frontier model
delivered as a managed cloud service: you call it over the internet, your data
goes to the vendor, you cannot see or own the weights, and it cannot run on a
disconnected device. That is the right product for many commercial uses and the
wrong one for a workload that cannot leave a controlled environment. Hanzo is the
opposite shape: open, owned-weight models you run on your own hardware---in a
data center, on a vehicle, or air-gapped---with post-quantum security built in,
your data never leaving, and cost driven down because you are not paying a
per-token markup or a hyperscaler margin. America's frontier AI has narrowed to
two closed labs; the only openly available frontier weights today are largely
foreign. For defense and regulated buyers who need open, sovereign, controllable
AI, that makes Hanzo the American alternative to both the closed duopoly and
foreign open weights.

\tableofcontents
\newpage

% ============================================================================
\section{The Real Axis of Comparison}
% ============================================================================

Discussions of ``which AI is better'' usually compare benchmark scores. For
commercial chat and coding that is a reasonable proxy. For defense and regulated
deployment it is the wrong axis. A model that scores two points higher on a
reasoning benchmark is of no use inside a controlled environment if it can only
be reached by sending the data outside that environment.

The decisive axis is the \emph{deployment and trust model}: who holds the weights,
where inference physically runs, whether the data leaves the enclave, whether the
system works without connectivity, what cryptography protects it, and what it
costs at scale. On this axis a closed cloud service and a sovereign self-hosted
stack are not competitors with different scores---they are different categories
of thing. Claude is the leading example of the first category; Hanzo is built as
the second.

We state plainly what closed cloud AI does well: it is operationally simple
(no infrastructure to run), it tracks the frontier of raw capability, and its
provider invests heavily in alignment and safety research. None of that is in
dispute. The argument here is narrow and structural: several requirements that
are mandatory for defense and regulated workloads cannot be met by a model you
do not hold and cannot run yourself.

% ============================================================================
\section{Dimension-by-Dimension}
% ============================================================================

\begin{table}[h]
\centering
\small
\begin{tabular}{p{0.27\textwidth} p{0.32\textwidth} p{0.32\textwidth}}
\hline
\textbf{Dimension} & \textbf{Closed cloud AI (e.g.\ Claude)} & \textbf{Hanzo (sovereign OSS)} \\
\hline
Model weights & Not disclosed; not transferable & Owned (Zen family, 0.8B--1T+); inspectable \\
Where inference runs & Vendor cloud only & Your hardware: data center, edge, air-gapped \\
Your data & Leaves the enclave to the vendor & Never leaves; stays where it lives \\
Disconnected / DDIL & Not possible (requires connectivity) & First-class (offline on embedded ARM) \\
Provenance & Vendor-controlled & American-owned, full provenance \\
Post-quantum security & Not a product property & Native (ML-KEM-768 + ML-DSA-65) \\
Auditability & Black box & Open stack; cryptographically signed actions \\
Customization & Prompting / limited fine-tune via API & Full fine-tune on your data, in your enclave \\
Cost model & Per-token metered & Owned + self-hosted; no markup or cloud margin \\
Lock-in & High (API + data gravity) & Low (you hold the stack) \\
\hline
\end{tabular}
\caption{Closed cloud AI versus a sovereign self-hostable stack. The comparison
is of deployment model, not of raw benchmark quality.}
\end{table}

\subsection{Data residency and the enclave boundary}
The single most consequential difference is the enclave boundary. To use a closed
cloud model, the data being reasoned over must cross the boundary to the vendor.
For CUI, engineering files, procurement records, or operational data this is often
simply prohibited. Hanzo runs the model inside the boundary, so the data never
crosses it. This is not a policy preference that can be waived with a contract
clause; it is a property of where computation physically happens.

\subsection{Disconnected and contested operation}
A cloud service is unavailable the moment the link drops or is jammed. At the
tactical edge, on a ship, or in an air-gapped facility there is no link to call.
Hanzo's models are quantized to run fully offline on hardware ranging from a
laptop to an embedded processor, so capability persists in exactly the
environments where a cloud service cannot operate at all.

\subsection{Provenance and the supply chain}
For an open, controllable alternative to a closed cloud model, today's practical
options are foreign open-weight models or Hanzo. Defense buyers care about who
produced the weights and whether they can be trusted and inspected. Hanzo's
position is specific: an American-owned, open-source frontier stack with owned
weights, so an organization can have open, auditable AI without depending on
either the closed domestic duopoly or foreign-origin weights.

\subsection{Post-quantum security}
Security against a future quantum adversary is a property of the whole stack, not
a feature of a chat API. Hanzo's transport, identity, and key management are
post-quantum by construction (all three NIST standards in production, formally
verified). A closed chat service does not expose this as something the customer
controls.

\subsection{Cost at scale}
A metered API charges per token on every inference, forever, with the vendor's
margin baked in. Owning the model and self-hosting removes both the per-token
markup and the hyperscaler margin; combined with an efficient transport layer
(the ZAP protocol's 91\% memory and 96\% energy reductions) this is the mechanism
behind roughly an order-of-magnitude lower infrastructure cost at scale. A
companion case study documents a regulated client whose cloud spend fell
approximately 97\% after migrating to this stack.

\subsection{Agentic throughput: one engineer, a team's output}
A decisive data point sits behind this whole comparison. Hanzo's founding CEO/CTO,
Zach Kelling, was---by public usage tracking (\texttt{claudecount.com})---the
\textbf{\#1 user of Anthropic by token count worldwide in 2025}. That throughput was
not spent on demos: on the flagship regulated production build, a single founding
engineer, multiplied by the agentic stack, authored the substantial majority of the
production code---thousands of commits in a matter of weeks---output that
conventionally requires a sizable team. The point for this paper is what came next.
That same throughput now runs on Hanzo's \emph{owned}, self-hostable stack rather
than a metered cloud API. The capability that topped a closed provider's global
usage is exactly the capability Hanzo packages as a sovereign, owned
alternative---so a buyer gets the productivity without the per-token meter, the data
egress, or the dependency. That is the comparison in one line: same frontier-grade
agentic output, but owned instead of rented.

% ============================================================================
\section{When Each Model Fits}
% ============================================================================

This is a decision framework, not a verdict. Closed cloud AI is the right choice
when the workload is non-sensitive, connectivity is reliable, operational
simplicity outweighs control, and there is no requirement to own or inspect the
model. A sovereign self-hosted stack is the right---often the only---choice when
any of the following hold: the data cannot leave a controlled environment; the
system must work disconnected or in contested conditions; provenance and
auditability are required; post-quantum protection is mandated; or per-token cost
at scale is untenable. Defense and regulated workloads concentrate in the second
column. That is the gap Hanzo is built to fill.

% ============================================================================
\section{Conclusion}
% ============================================================================

The honest comparison is not ``Hanzo's model beats Claude.'' It is that a closed
cloud model and a sovereign self-hostable stack are different categories, and the
requirements of defense and regulated workloads---data residency, disconnected
operation, provenance, post-quantum security, auditability, and cost---select the
second category by construction. Within that category the choice narrows to
foreign open weights or an American open-source alternative. Hanzo is built to be
that alternative: own the model, own the cryptography, run it where the data
lives.

\end{document}
